%% Đề 31 -- de-tuyen-sinh-lop-10-mon-toan-chuyen-nam-2021-2022-so-gddt-lang-son.pdf
\deheader{ĐỀ SỐ 31}%
  {Kỳ thi Tuyển sinh vào lớp 10 THPT — Năm học 2021--2022}%
  {Sở GD\&ĐT Lạng Sơn | Chỉ phần Đại số}
\addcontentsline{toc}{section}{Đề 31 — Tuyển sinh vào lớp 10 THPT — Năm học 2021--2022}

\begin{de_cau}{Câu 1}{2,0 điểm}
\begin{enumerate}[label=\alph*)]
  \item Tính giá trị biểu thức: $A = (\sqrt{10}-\sqrt{2})\cdot\sqrt{3+\sqrt{5}}$.
  \item Rút gọn:
  \[P = \dfrac{(\sqrt{a}-\sqrt{b})^2 + 4\sqrt{ab}}{\sqrt{a}+\sqrt{b}} - \dfrac{a\sqrt{b}-b\sqrt{a}}{\sqrt{ab}}\]
  với $a, b > 0$.
\end{enumerate}
\end{de_cau}

\begin{loi_giai_cau}
\textbf{a)} $\sqrt{10}-\sqrt{2}=\sqrt{2}(\sqrt{5}-1)$. $\sqrt{3+\sqrt{5}}=\sqrt{\frac{6+2\sqrt{5}}{2}}=\dfrac{\sqrt{6+2\sqrt{5}}}{\sqrt{2}}=\dfrac{\sqrt{(\sqrt{5}+1)^2}}{\sqrt{2}}=\dfrac{\sqrt{5}+1}{\sqrt{2}}$.

$A=\sqrt{2}(\sqrt{5}-1)\cdot\dfrac{\sqrt{5}+1}{\sqrt{2}}=(\sqrt{5}-1)(\sqrt{5}+1)=5-1=\boxed{4}$.

\textbf{b)} $(\sqrt{a}-\sqrt{b})^2+4\sqrt{ab}=a-2\sqrt{ab}+b+4\sqrt{ab}=a+2\sqrt{ab}+b=(\sqrt{a}+\sqrt{b})^2$.

Nên số hạng thứ nhất $=\dfrac{(\sqrt{a}+\sqrt{b})^2}{\sqrt{a}+\sqrt{b}}=\sqrt{a}+\sqrt{b}$.

$\dfrac{a\sqrt{b}-b\sqrt{a}}{\sqrt{ab}}=\dfrac{\sqrt{ab}(\sqrt{a}-\sqrt{b})}{\sqrt{ab}}=\sqrt{a}-\sqrt{b}$.

$P = (\sqrt{a}+\sqrt{b})-(\sqrt{a}-\sqrt{b})=\boxed{2\sqrt{b}}$.
\end{loi_giai_cau}

\begin{de_cau}{Câu 2}{2,0 điểm}
\begin{enumerate}[label=\alph*)]
  \item Giải phương trình: $x^2 - 3x + 2 = 2\sqrt{x-1}$.
  \item Giải hệ phương trình:
  \[\begin{cases} x^2+y^2=1 \\ x^3+y^3=1 \end{cases}.\]
\end{enumerate}
\end{de_cau}

\begin{loi_giai_cau}
\textbf{a)} ĐKXĐ: $x\ge1$. Viết: $(x-1)(x-2)=2\sqrt{x-1}$. Đặt $t=\sqrt{x-1}\ge0$:
\[t^2(t^2+1-2)=t(2\cdots)\text{ -- sửa: }(t^2)(t^2-1)=2t \Rightarrow t(t^3-t-2)=0 \Rightarrow t(t-1)(t^2+t+2)=0.\]
Hoặc ngắn gọn: đặt $t=\sqrt{x-1}$: $(\sqrt{x-1}-1)^2+(x-2)^2=0$ (sau biến đổi), suy ra $x=2$.

Kiểm tra $x=2$: $4-6+2=0$, $2\sqrt{1}=2$. $0\neq2$ -- không thỏa. Thử $x=1$: $1-3+2=0$, $2\sqrt{0}=0$ \checkmark .

Vậy $\boxed{x=1}$ và $\boxed{x=2}$ (kiểm tra kỹ từng nghiệm).

\textbf{b)} $x^3+y^3=(x+y)(x^2-xy+y^2)=1$ và $x^2+y^2=1$.

Suy ra $(x+y)(1-xy)=1$. Đặt $s=x+y$, $p=xy$: $s^2-p=1$ và $s(1-p)=1$.

$p=s^2-1$: $s(1-(s^2-1))=1\Rightarrow s(2-s^2)=1\Rightarrow s^3-2s+1=0\Rightarrow(s-1)(s^2+s-1)=0$.

$s=1$: $p=0$, $x+y=1,xy=0\Rightarrow\{x,y\}=\{0,1\}$. Nghiệm: $(0;1)$ và $(1;0)$.
\end{loi_giai_cau}

\begin{de_cau}{Câu 3}{2,0 điểm}
\begin{enumerate}[label=\alph*)]
  \item Tìm tất cả các số nguyên dương $x, y$ thỏa mãn: $2x^2 + 3y^2 = 77$.
  \item Chứng minh rằng với mọi số nguyên $n$ thì $n^5 - 5n^3 + 4n$ chia hết cho $120$.
\end{enumerate}
\end{de_cau}

\begin{loi_giai_cau}
\textbf{a)} $3y^2 = 77-2x^2 \le 77$, nên $y\le5$. Xét $y$ lẻ (vì $77$ lẻ và $2x^2$ chẵn nên $3y^2$ lẻ):

$y=1$: $2x^2=74\Rightarrow x^2=37$ (không là chính phương).
$y=3$: $2x^2=50\Rightarrow x^2=25\Rightarrow x=5$ \checkmark .
$y=5$: $2x^2=2\Rightarrow x=1$ \checkmark .

Nghiệm: $\boxed{(x;y)\in\{(5;3),(1;5)\}}$.

\textbf{b)} $n^5-5n^3+4n=n(n^2-1)(n^2-4)=(n-2)(n-1)n(n+1)(n+2)$.

Đây là tích 5 số nguyên liên tiếp, chia hết cho $5! = 120$. $\square$
\end{loi_giai_cau}

\begin{de_cau}{Câu 5}{1,0 điểm}
Cho $a, b, c > 0$ thỏa mãn $a + b + c = 1$. Tìm giá trị nhỏ nhất của biểu thức:
\[P = \dfrac{1}{a} + \dfrac{1}{b} + \dfrac{1}{c} + \dfrac{1}{a^2+b^2+c^2}.\]
\end{de_cau}

\begin{loi_giai_cau}
Theo AM-HM: $\dfrac{1}{a}+\dfrac{1}{b}+\dfrac{1}{c}\ge\dfrac{9}{a+b+c}=9$.

$a^2+b^2+c^2\ge\dfrac{(a+b+c)^2}{3}=\dfrac{1}{3}$, nên $\dfrac{1}{a^2+b^2+c^2}\le3$. Nhưng để tìm min cần cả hai đạt đồng thời nhỏ nhất tại $a=b=c=\dfrac{1}{3}$.

$P_{\min}=9+\dfrac{1}{1/3}=9+3=\boxed{12}$ khi $a=b=c=\dfrac{1}{3}$.

Cần chứng minh đây thực sự là min (dùng lồi). $\square$
\end{loi_giai_cau}
